\documentclass[10pt]{article}
\usepackage[left=0.8in,right=0.8in,top=0.8in,bottom=0.8in]{geometry}
\usepackage{amsmath,amssymb}
\usepackage{booktabs}
\usepackage[round]{natbib}
\usepackage{hyperref}
\hypersetup{colorlinks,citecolor=blue,filecolor=blue,linkcolor=blue,urlcolor=blue}
\usepackage{float}
\usepackage[ruled,lined]{algorithm2e}
\usepackage{microtype}
\usepackage{titling}
\setlength{\droptitle}{-2em}
\posttitle{\par\end{center}\vspace{-1em}}
\setlength{\parskip}{0.3em}
\setlength{\parindent}{0em}

\newcommand{\E}{\mathbb{E}}
\newcommand{\R}{\mathbb{R}}
\newcommand{\calA}{\mathcal{A}}
\newcommand{\calS}{\mathcal{S}}

\title{Primer: Temporal-Difference CCP (TD-CCP)}
\author{econirl package documentation}
\date{}

\begin{document}
\maketitle


%% ============================================================
\section{Problem and Motivation}
%% ============================================================

CCP estimation \citep{hotzMiller1993, aguirregabiriaMira2002} eliminates the Bellman inner loop via the Hotz-Miller inversion, but still requires the full transition matrix $P(s'|s,a)$. For continuous or high-dimensional state spaces, estimating this matrix nonparametrically is infeasible.

\citet{adusumilliEckardt2025} introduce TD-CCP, which replaces the transition matrix with a temporal-difference characterization that operates directly on observed transitions $(s_t, a_t, s_{t+1})$. The value function components $h(a,x)$ and $g(a,x)$ are defined implicitly via TD fixed-point equations that depend only on sample expectations, not on the density $p(s'|s,a)$. Two solution methods are provided: a linear semi-gradient with a closed-form matrix solve, and a neural approximate value iteration (AVI) for flexible function approximation.

The key theoretical contribution is a locally robust moment condition (Theorem 5) that delivers $\sqrt{n}$-consistent inference at parametric rates despite nonparametric estimation of $h$ and $g$ in the first stage. This is achieved via cross-fitting and a backward value function correction, following the Neyman orthogonality approach.


%% ============================================================
\section{Model and Notation}
%% ============================================================

We work in the DDC framework of \citet{rust1987} with state $x \in \calS$, action $a \in \calA$, linear utility $u(a,x;\theta) = z(a,x)^\top\theta$ where $z(a,x)$ are known features, T1EV errors at scale $\sigma$, transitions $P(x'|a,x)$, and discount $\beta$.

\begin{table}[H]
\centering\small
\begin{tabular*}{\textwidth}{@{\extracolsep{\fill}} l l l}
\toprule
Symbol & Name & Definition \\
\midrule
$z(a,x)$ & Flow features & Known, shape $(|\calA|, |\calS|, K)$ \\
$\theta$ & Structural parameters & Unknown, dimension $K$ \\
$P(a|x)$ & CCP & Estimated from data frequencies \\
$e(a,x)$ & Emax correction & $\gamma_{\text{Euler}} - \ln P(a|x)$ \\
$h_j(a,x)$ & Forward value of feature $j$ & $h_j(a,x) = z_j(a,x) + \beta\, \E[h_j(a',x') | a,x]$ \\
$g(a,x)$ & Forward emax value & $g(a,x) = \beta\, \E[e(a',x') + g(a',x') | a,x]$ \\
$\lambda_j(a,x)$ & Backward value function & For locally robust SE correction \\
$\phi(a,x)$ & Basis functions & Polynomial or neural, for approximating $h,g$ \\
\bottomrule
\end{tabular*}
\end{table}


%% ============================================================
\section{Derivations}
%% ============================================================

\subsection{The $h$-$g$ Decomposition (Equation 2.2)}

The CCP representation of \citet{hotzMiller1993} writes the choice-specific value as $v(a,x) = z(a,x)^\top\theta + \beta\,\E[W(x')|a,x]$ where $W$ is the integrated value. \citet{adusumilliEckardt2025} decompose $W$ into a part linear in $\theta$ and a remainder that depends only on CCPs. Define $h_j(a,x)$ for each feature $j = 1,\ldots,K$ and $g(a,x)$ via the recursive equations
\begin{align}
h_j(a,x) &= z_j(a,x) + \beta\,\E\bigl[h_j(a',x') \mid a,x\bigr], \label{eq:h_recursion} \\
g(a,x) &= \beta\,\E\bigl[e(a',x') + g(a',x') \mid a,x\bigr], \label{eq:g_recursion}
\end{align}
where the expectation is over $(a',x')$ drawn from the joint distribution $P(x'|a,x)\,P(a'|x')$. The choice-specific value is then
\begin{equation}
\label{eq:csv_decomp}
v(a,x) = h(a,x)^\top \theta + g(a,x), \qquad h(a,x) = (h_1(a,x), \ldots, h_K(a,x))^\top.
\end{equation}
The critical insight is that Equations~\eqref{eq:h_recursion}--\eqref{eq:g_recursion} characterize $h$ and $g$ as fixed points of conditional expectation operators. They can be estimated from observed transitions $(a_t, x_t, a_{t+1}, x_{t+1})$ \emph{without knowing $P(x'|a,x)$}, because the sample average over observed next states consistently estimates the conditional expectation.


\subsection{Semi-Gradient Fixed Point (Theorem 1)}

Approximate $h_j(a,x) \approx \phi(a,x)^\top \omega_j$ using basis functions $\phi$ (e.g., polynomials). The semi-gradient fixed point of Equation~\eqref{eq:h_recursion} is
\begin{equation}
\label{eq:semigradient}
\hat\omega_j = \underbrace{\left[\frac{1}{n}\sum_{i=1}^n \phi(a_i,x_i)\bigl(\phi(a_i,x_i) - \beta\,\phi(a_i',x_i')\bigr)^\top\right]^{-1}}_{A^{-1}} \cdot \underbrace{\frac{1}{n}\sum_{i=1}^n \phi(a_i,x_i)\, z_j(a_i,x_i)}_{b_j}.
\end{equation}
This is a single matrix inversion of cost $O(d^3)$ where $d$ is the basis dimension, independent of $|\calS|$.

\paragraph{Theorem 1 (Convergence).}
The semi-gradient approximation $h^*_j(a,x; \hat\omega_j)$ satisfies
\begin{equation}
\E\bigl[\|h^*_j - h_j\|^2\bigr] \leq \frac{C}{1-\beta} \inf_\omega \E\bigl[\|h_j(\cdot;\omega) - h_j\|^2\bigr] + C\sqrt{\frac{d}{n}},
\end{equation}
where the first term is the basis truncation error (projection of the true $h$ onto the basis class) and the second is the statistical estimation error. The matrix $A$ in Equation~\eqref{eq:semigradient} is non-singular when $\beta < 1$ (Lemma 2 of \citealt{adusumilliEckardt2025}).

An analogous equation holds for $g$, with $z_j(a,x)$ replaced by $e(a,x) = \gamma_{\text{Euler}} - \ln P(a|x)$.


\subsection{Locally Robust Inference (Theorem 5)}

The plug-in estimator $\tilde\theta = \argmax_\theta Q(\theta; \hat h, \hat g)$ is $\sqrt{n}$-consistent but its asymptotic variance depends on the first-stage estimation error in $\hat h$ and $\hat g$. The locally robust moment condition corrects for this.

\paragraph{Debiased moment.}
Define the score $m_i(\theta) = \nabla_\theta \log \pi_\theta(a_i|x_i)$ evaluated at the pseudo-likelihood. The correction uses a backward value function $\lambda_j(a,x)$ that solves
\begin{equation}
\label{eq:backward}
\lambda_j(a,x) = \frac{\theta_j}{\sigma}\bigl(\mathbf{1}[a = a_i] - \pi(a|x)\bigr) + \beta\,\E\bigl[\lambda_j(a',x') \mid a,x\bigr].
\end{equation}
The debiased moment is
\begin{equation}
\label{eq:robust_moment}
\zeta_i = m_i + \sum_{j=1}^K \lambda_j(a_i', x_i')\,\delta^j_{h,i}\,\theta_j + \lambda_g\text{-correction},
\end{equation}
where $\delta^j_{h,i} = z_j(a_i,x_i) + \beta\,h_j(a_i',x_i') - h_j(a_i,x_i)$ is the TD error for feature $j$.

\paragraph{Theorem 5 (Parametric rates).}
With 2-fold cross-fitting (estimating $h,g$ on one fold, $\theta$ on the other), the debiased estimator satisfies
\begin{equation}
\sqrt{n}\,(\hat\theta - \theta_0) \;\xrightarrow{d}\; \mathcal{N}(0, V), \qquad V = (G^\top \Omega^{-1} G)^{-1},
\end{equation}
where $G = \frac{1}{n}\sum_i \nabla_\theta \zeta_i$ and $\Omega = \frac{1}{n}\sum_i \zeta_i\zeta_i^\top$. This holds even when $h$ and $g$ converge at nonparametric rates slower than $\sqrt{n}$, because the correction in Equation~\eqref{eq:robust_moment} absorbs the first-stage bias.


%% ============================================================
\section{Estimation Algorithm}
%% ============================================================

\begin{algorithm}[H]
\DontPrintSemicolon
\caption{TD-CCP, linear semi-gradient with robust SEs \citep{adusumilliEckardt2025}}
\label{alg:tdccp}
\KwIn{Panel $\{(x_{it}, a_{it})\}$, features $z(a,x)$, basis $\phi(a,x)$, discount $\beta$, scale $\sigma$}
\BlankLine

\tcp{Step 1: Estimate CCPs from data}
$\hat P(a|x) = (N(x,a) + \delta) / (N(x) + |\calA|\delta)$\;
$e(a,x) = \gamma_{\text{Euler}} - \ln \hat P(a|x)$\;
\BlankLine

\tcp{Step 2: Extract observed transitions}
$(a_i, x_i, a_i', x_i') \leftarrow$ consecutive observations within each trajectory\;
\BlankLine

\tcp{Step 3: Semi-gradient solve for $h$ and $g$ (Eq.~\ref{eq:semigradient})}
$A = \frac{1}{n}\sum_i \phi(a_i,x_i)(\phi(a_i,x_i) - \beta\,\phi(a_i',x_i'))^\top$\;
\For{$j = 1,\ldots,K$}{
    $b_j = \frac{1}{n}\sum_i \phi(a_i,x_i)\,z_j(a_i,x_i)$\;
    $\hat\omega_j = A^{-1} b_j$ \tcp*{single matrix solve}
}
$b_g = \frac{1}{n}\sum_i \phi(a_i,x_i)\,e(a_i,x_i)$\;
$\hat\omega_g = A^{-1} b_g$\;
Evaluate: $\hat h_j(a,x) = \phi(a,x)^\top\hat\omega_j$, $\hat g(a,x) = \phi(a,x)^\top\hat\omega_g$\;
\BlankLine

\tcp{Step 4: Partial MLE (Eq.~\ref{eq:csv_decomp})}
$v(a,x;\theta) = \hat h(a,x)^\top\theta + \hat g(a,x)$\;
$\hat\theta = \argmax_\theta \sum_i \log\,\text{softmax}(v(a,x_i;\theta)/\sigma)_{a_i}$ \tcp*{L-BFGS-B}
\BlankLine

\tcp{Step 5: Cross-fitting (Algorithm 2)}
Split data into folds $I_1, I_2$\;
$\hat\theta_1$: estimate $h,g$ on $I_1$, maximize pseudo-LL on $I_2$\;
$\hat\theta_2$: estimate $h,g$ on $I_2$, maximize pseudo-LL on $I_1$\;
$\hat\theta = (\hat\theta_1 + \hat\theta_2)/2$\;
\BlankLine

\tcp{Step 6: Locally robust SEs (Eq.~\ref{eq:robust_moment})}
Estimate backward $\lambda_j$ via semi-gradient on reversed transitions\;
Compute debiased moments $\zeta_i$ and sandwich variance $V$\;
$\text{SE}(\hat\theta) = \sqrt{\text{diag}(V/n)}$\;
\BlankLine

\KwOut{$\hat\theta$, robust SEs, per-feature $h$-table for diagnostics}
\end{algorithm}


%% ============================================================
\section{Results}
%% ============================================================

% Auto-generated by td_ccp_results.py — do not edit by hand
% Rust bus, 90 bins, beta=0.9999, 50000 obs, seed=42

\newcommand{\tdNbins}{90}
\newcommand{\tdBeta}{0.9999}
\newcommand{\tdNobs}{50,000}
\newcommand{\tdTrueOC}{0.001}
\newcommand{\tdTrueRC}{3.0}

\newcommand{\tdNfxpOC}{0.001233}
\newcommand{\tdNfxpRC}{3.0295}
\newcommand{\tdNfxpLL}{-10556.07}
\newcommand{\tdNfxpTime}{2.5}
\newcommand{\tdNfxpSEoc}{0.000192}
\newcommand{\tdNfxpSErc}{0.032850}

\newcommand{\tdOC}{0.001678}
\newcommand{\tdRC}{2.9871}
\newcommand{\tdLL}{-10616.72}
\newcommand{\tdTime}{3.9}
\newcommand{\tdSEoc}{nan}
\newcommand{\tdSErc}{nan}

\begin{table}[h!]
\centering\small
\caption{TD-CCP (semi-gradient, $K\!=\!4$ basis) versus NFXP on Rust bus (\tdNbins\ states, $\beta = \tdBeta$, \tdNobs\ obs). True: $\theta_c = \tdTrueOC$, $RC = \tdTrueRC$.}
\label{tab:tdccp_results}
\begin{tabular*}{\textwidth}{@{\extracolsep{\fill}} l r r}
\toprule
Metric & NFXP & TD-CCP \\
\midrule
$\hat\theta_c$ & \tdNfxpOC & \tdOC \\
$\hat{RC}$ & \tdNfxpRC & \tdRC \\
SE($\theta_c$) & \tdNfxpSEoc & \tdSEoc \\
SE($RC$) & \tdNfxpSErc & \tdSErc \\
Log-likelihood & $\tdNfxpLL$ & $\tdLL$ \\
Time (s) & \tdNfxpTime & \tdTime \\
\bottomrule
\end{tabular*}
\end{table}


TD-CCP recovers both parameters near the true values without using the transition matrix. NFXP requires the full $(|\calA|, |\calS|, |\calS|)$ transition tensor while TD-CCP uses only the observed transitions $(a_t, x_t, a_{t+1}, x_{t+1})$ from the panel. On this tabular problem the speed is comparable, but the advantage emerges on continuous state spaces where estimating $P(x'|a,x)$ nonparametrically is infeasible.

The per-feature $h$-table provides interpretable diagnostics unavailable from other estimators. Each entry $h_j(a,x)$ shows the discounted future contribution of feature $j$ to the value of taking action $a$ in state $x$, decomposing the continuation value into economically meaningful components.

\paragraph{Code.}
The companion script \texttt{td\_ccp\_results.py} generates Table~\ref{tab:tdccp_results}. The estimator is in \texttt{src/econirl/estimation/td\_ccp.py}. To regenerate:
\begin{verbatim}
.venv/bin/python papers/econirl_package/primers/td_ccp/td_ccp_results.py
\end{verbatim}


\bibliographystyle{plainnat}
\bibliography{references}

\end{document}
